% This is a file that contains the tracking of the activities related to the project during the first week of March 2025.
% The activities are divided into groups and a summary of the progress is provided.

\documentclass{article}

\usepackage{graphicx}
\usepackage[backend=biber,style=numeric]{biblatex}

\newcommand{\apj}{The Astrophysical Journal}

\addbibresource{references.bib}

\begin{document}
\author{Joan Ronquillo}

\title{Project Progress - Week 1 of March 2025}
\maketitle

This is a file that contains the tracking of the activities related to the OFC project during the first week of March 2025.
The activities are divided into groups and a summary of the progress is provided.

\section{Initial Status of the Project}
The project has made significant progress up until the first week of March 2025. The key achievements include:

\begin{itemize}
    \item Successful implementation and validation of the optical force calculation function in the dipole approximation.
    \item Development of test and validation scripts that confirm the correct dependence on the permittivity.
    \item Advancement in the theoretical documentation, particularly in the section on polarizability.
    \item Detailed derivation of the quasiestatic polarizability formula for a dielectric sphere.
\end{itemize}

\section{Potential Tasks for the Week}
The following tasks have been identified as potential areas of focus for the first week of March 2025:
\begin{itemize}
    \item Optimize the functions for querying the refractiveindex.info database.
    \item Complete the integration of permittivities into an specific equation.
    \item Finish documenting and providing examples of usage.
\end{itemize}

\section{Week Tracking}
\subsection{Monday - March 3, 2025}
The repository structure has been fixed and the management of Python modules has been addressed.
Additionally, I have been learning how to compile LaTeX projects located inside the repository.
This document serves as an example of the compilation process.

\subsection{Tuesday - March 4, 2025}
The tasks for the week have been specified and the progress in documentation, code, and testing has been tracked.
We are going to focus on the algorithm used to obtain the polarizability of the particles. In the "dipole momenta" branch,
a first version of ".h" and ".cpp" files has been created as a first approach to modular development of the functionalities.
I have encountered some issues with warnings and errors flagged by the editor in VSCode. While the code works properly, the editor
flags some issues that are not real. It may be related to the editor's settings or installation.
Apart from this, it has been observed that a CMakefile must be created, an include directory must be added,
and we need to explore how to include different functions without having to import multiple files.

\subsection{Thursday - March 6, 2025}
The reading of Draine's paper on radiative corrections in polarizability has been initiated.
In addition, work has been done on including BibTeX references in LaTeX files.
A references file has been created and included in the main file. In this regard,
the directory structure of the tracking has been modified.
A directory tree has been added to the document to visualize the structure of the tracking directory.

\section{Current Next Steps}
The next steps for the project include:
\begin{itemize}
    \item Explor de CMakefile, include directory, and modular development (Notion).
    \item Development of a calculation method for Optical Forces given a polarizability, a field, and a field gradient.
    \item Understand radiative correction in polarizability. Check the paper \cite{Draine1988ApJ...333..848D}.
    \item Documentation of quasiestatic polarizability theory. 
\end{itemize}

\printbibliography

\end{document}


